\chapter{Introduction}

%------------------------------- 1.1 -------------------------------

\section{Background}

Sensitive information is increasingly present in digital visual media. Photographs and images are widely used across domains such as journalism, surveillance, documentation, mapping, and social platforms. While visual data offers significant practical value, it also introduces challenges related to privacy, security, and ethical responsibility. Images may contain sensitive information, including faces, license plates, textual content, and contextual details that can directly or indirectly identify individuals \cite{frome2009large}.

When such information is disclosed without proper handling, it may lead to privacy violations, unauthorized identification, or misuse of personal data. As a consequence, regulations such as the General Data Protection Regulation (GDPR) impose strict requirements on how personal data must be processed and protected \cite{gdpr2016}. Organizations and individuals handling visual data must ensure that sensitive content is properly anonymized prior to publication, storage, or sharing.

Traditionally, anonymizing sensitive information in images has relied on manual editing using image processing software. While manual redaction provides a high degree of control, it is inherently time-consuming, repetitive, and prone to human error. These limitations become particularly significant when processing large volumes of visual material, where manual approaches become inefficient and difficult to scale \cite{piao2024deidentification}.

Recent advances in artificial intelligence and machine learning have enabled automated techniques for detecting sensitive content in images. Modern computer vision models, particularly those based on deep learning, can identify objects such as faces, license plates, and textual elements with increasing accuracy \cite{raina2023egoblur}. These developments create opportunities for more efficient anonymization workflows that reduce manual effort while improving consistency. However, existing solutions still vary considerably in terms of detection accuracy, usability, and adaptability to different application contexts.

%------------------------------- 1.2 -------------------------------

\section{Problem Statement}

Despite the availability of computer vision techniques for detecting sensitive content, many existing solutions remain limited in terms of accuracy, usability, and privacy considerations. Studies on privacy-preserving computer vision systems highlight challenges related to reliable detection performance and the difficulty of achieving consistent anonymization across diverse visual environments \cite{piao2024deidentification}. In addition, some solutions rely on cloud-based processing or centralized storage, which may introduce additional privacy concerns and complicate integration into secure organizational workflows \cite{raina2023egoblur}. These limitations make it challenging for organizations to adopt automated solutions that both protect privacy and fit into existing operational requirements.

Safe Media AI AS has developed an AI-based platform primarily focused on text-based document processing. As visual data continues to grow in importance, extending automated redaction capabilities to images represents a natural and necessary progression. However, designing an effective system for automated image redaction requires addressing challenges related to detection accuracy, usability, and privacy-by-design principles.

This project investigates how a system, named SafeVision, can be designed and implemented to automatically identify and anonymize sensitive information in still images while maintaining efficiency, usability, and strong privacy protection. In addition to reliable detection and redaction, SafeVision should provide flexible user interaction and interface design that allows users to review, adjust, and control the anonymization process according to their specific workflow requirements.

%------------------------------- 1.3 -------------------------------

\section{Project Goals and Scope}

The main goal of this thesis is to design and develop a prototype system capable of automatically detecting and anonymizing sensitive information in still images, and later extended to video. The system aims to identify elements such as faces, license plates, and textual content and apply appropriate redaction techniques to prevent unintended disclosure of personal or confidential information.

In addition to automated detection, the solution will provide interactive editing capabilities that allow users to review and adjust redaction results through a graphical user interface. This approach combines automated computer vision techniques with human oversight to improve both reliability and usability.

The scope of the project is limited to the development and evaluation of a functional prototype for still image processing. The system will support common image formats such as PNG and JPEG and will be designed to integrate with the existing platform developed by Safe Media AI AS. This means, initial support for PDF is no longer considered. The solution will also follow privacy-by-design principles and consider relevant data protection requirements, including GDPR.

%------------------------------- 1.4 -------------------------------

\section{Delimitation}

The scope of this project is limited to the development of a prototype system for identifying and anonymizing sensitive information in visual media, including both still images and video. The primary focus is on establishing a reliable detection and redaction pipeline rather than developing a fully optimized large-scale processing system.

The solution is designed to integrate with the existing platform developed by Safe Media AI AS, while also supporting standalone use. However, extensive customization for different third-party systems or deployment environments is outside the scope of this project.

Infrastructure constraints also influence the implementation. The system is developed and tested using Google Cloud services under a free-tier environment, which limits computational resources and requires model inference to run on CPU rather than GPU. In addition, only machine learning models with licenses compatible with commercial use are considered. Models with restrictive licenses that prevent commercial deployment are excluded from the project.

%------------------------------- 1.5 -------------------------------

\section{Significance and Impact}

The increasing reliance on visual media in digital systems creates a growing need for effective privacy-preserving technologies. Images and video frequently contain identifiable information such as faces, license plates, or textual elements, which must be carefully handled to comply with data protection regulations such as the General Data Protection Regulation (GDPR) \cite{gdpr2016}. As organizations process larger volumes of visual data, manual redaction becomes increasingly inefficient and difficult to scale \cite{piao2024deidentification}.

This project contributes to the development of automated anonymization techniques by exploring the use of computer vision for detecting and redacting sensitive content in images and video. The proposed system combines automated detection with interactive editing features, enabling users to review and adjust anonymization results. In addition to its academic relevance within privacy-preserving computer vision, the project supports Safe Media AI AS in extending its platform from text-based document processing to visual media, providing a foundation for future development of automated privacy-preserving tools.

%------------------------------- 1.6 -------------------------------

\section{Project Report Structure}

This thesis is structured to reflect the progression of the project, from problem formulation and theoretical foundation to system design, implementation, evaluation, discussion, and conclusion.

The introductory chapters establish the context and motivation for the project. Chapter 1 presents the background, problem statement, project goals, scope, delimitations, chosen approach, and significance of the work. Chapter 2 provides the theoretical and technical foundation by discussing sensitive information in visual media, GDPR-related design constraints, existing redaction workflows, AI-based redaction methods, and the model choices used in SafeVision.

Following this foundation, the thesis moves into requirements, process, and design. Chapter 3 describes the development process, including the project methodology, sprint structure, stakeholder involvement, project timeline, and working practices. Chapter 4 defines the functional and non-functional requirements for the system, including use cases and requirement prioritization. Chapter 5 presents the system architecture, including the overall architectural design, main system components, machine learning pipeline, and deployment architecture.

The next chapters describe the design and implementation of SafeVision. Chapter 6 focuses on the user interface and usability, explaining the design methodology, interface components, and interaction choices that support user review and control (STANDALONE). Chapter 7 presents the technologies and methodologies used in the project, including frontend, backend, machine learning, storage, deployment, and infrastructure choices. Chapter 8 describes the implementation details of the frontend, backend, AI/ML integration, redaction pipeline, video processing, and deployment setup.

The later chapters evaluate and reflect on the completed system. Chapter 9 presents the testing and system evaluation, including frontend testing, backend testing, user testing, performance testing, and security evaluation. Chapter 10 analyzes the sustainability of the system across human, social, environmental, economic, and technical dimensions. Chapter 11 presents the results from the pre-trained models, custom-trained detection models, image segmentation, and video processing.

Finally, Chapter 12 discusses the project outcomes, including achieved goals, design decisions, technical limitations, ethical and privacy considerations, adoption considerations, and future development opportunities. Chapter 13 concludes the thesis by summarizing the main contributions and findings. The appendices provide additional supporting material, including extended interface design, detailed model evaluation, training configurations, test coverage, metadata handling, the project plan, and selected meeting notes.